% Traditional advisor-facing RA-L paper outline for ToolFrame 6.4.
\documentclass[journal]{IEEEtran}

\usepackage{amsmath,amssymb}
\usepackage{booktabs}
\usepackage{xcolor}

\newcommand{\todo}[1]{\textcolor{red}{TODO: #1}}
\newcommand{\method}{ToolFrame}

\title{\method{}: Let the Tool Define the Frame for Dexterous Manipulation}
\author{Author~Names}

\begin{document}
\maketitle

\section*{Title Candidates}
\begin{enumerate}
  \item \method{}: Let the Tool Define the Frame for Dexterous Manipulation
  \item \method{}: Tool-First Visual Grounding for Dexterous Manipulation
  \item \method{}: Tool-Centric Visual Grounding for Dexterous Tool Use
  \item \method{}: Compiling Tool-Centric Visual Evidence into Dexterous Manipulation Subgoals
  \item Planning Dexterous Tool Use Around Functional Tool Geometry
\end{enumerate}

\section*{Core Thesis}
\method{} makes the tool, rather than the hand, the organizing frame for
dexterous manipulation. The VLM selects sparse public scene references and
functional tool references; a geometry-constrained compiler turns them into
whole-tool $SE(3)$ subgoals; a SimToolReal fixed-goal policy makes the hand
follow those tool-centric goals; task contracts evaluate final success and
attribute failures.

\section*{Abstract Skeleton}
\begin{enumerate}
  \item Problem: dexterous tool use is defined by the tool: which part contacts
  the scene, which axis aligns with the action, and where the whole tool should
  move.
  \item Gap: directly asking VLMs for robot trajectories entangles semantic
  grounding, metric geometry, controller feasibility, and hidden benchmark
  assumptions.
  \item Method: \method{} lets the tool define the planning frame by restricting
  VLM outputs to sparse public references: tool.contact\_point,
  tool.motion\_axis, start/end/path scene refs, and scene.direction.
  \item Compiler: lifted RGB-D scene refs and tool refs are compiled into
  whole-tool $SE(3)$ subgoals.
  \item Evaluation: SimToolReal fixed-goal rollout is scored with task
  contracts and module-level failure attribution.
  \item Results: \todo{fill final N, task success rates, and failure
  attribution only from frozen launcher\_summary files.}
\end{enumerate}

\section{Introduction}
\subsection{Motivation}
\begin{itemize}
  \item Dexterous tool use is tool-centric: the task is defined by functional
  tool geometry before it is defined by hand motion.
  \item The robot must know which tool part should contact the scene and how
  the whole tool should be oriented.
  \item Examples: brush sweep, eraser wipe, hammer strike, screwdriver insert,
  marker trace.
\end{itemize}

\subsection{Problem}
\begin{itemize}
  \item VLMs can identify task-relevant visual evidence but should not be
  trusted to produce metric robot trajectories.
  \item Trajectory-style VLM planning can hide unsupported target axes, contact
  locations, task-family defaults, or simulator state.
\end{itemize}

\subsection{Key Idea}
\begin{itemize}
  \item Use VLMs only for sparse, auditable tool and scene references.
  \item Use RGB-D lifting and geometric compilation for metric whole-tool
  subgoals.
  \item Let the hand follow tool-centric subgoals instead of asking the VLM to
  generate hand trajectories.
  \item Use task contracts for final success and failure attribution.
\end{itemize}

\subsection{Contributions}
\begin{enumerate}
  \item An auditable VLM interface for tool use: sparse scene refs and
  functional tool refs, not robot trajectories.
  \item A geometry-constrained compiler that binds tool.contact\_point and
  tool.motion\_axis to lifted RGB-D scene refs and emits whole-tool $SE(3)$
  subgoals.
  \item A task-contract benchmark protocol that separates headline tasks from
  candidate, diagnostic, and draft tasks, and reports module-level failure
  attribution.
\end{enumerate}

\section{Related Work}
\subsection{VLMs and Language-Conditioned Robot Planning}
\begin{itemize}
  \item LLM/VLM planners, code-as-policies systems, image-conditioned action
  proposal, affordance prompting.
  \item Positioning: \method{} uses VLMs for reference grounding, not trajectory
  generation.
\end{itemize}

\subsection{Tool-Use Affordances and Functional Part Grounding}
\begin{itemize}
  \item Tool affordance learning, part segmentation, contact-point prediction,
  functional region prediction.
  \item Positioning: \method{} converts public tool-part evidence into
  tool.contact\_point and tool.motion\_axis.
\end{itemize}

\subsection{Constraint-Based Manipulation and Subgoal Planning}
\begin{itemize}
  \item Task-space constraints, contact-aware manipulation, keyframe/subgoal
  planning, $SE(3)$ optimization.
  \item Positioning: sweep/strike/wipe/insert are task language, not compiler
  primitives.
\end{itemize}

\subsection{Robotic Manipulation Benchmarks}
\begin{itemize}
  \item Manipulation benchmarks, tool-use benchmarks, success metrics.
  \item Positioning: \method{} uses task contracts to report final success and
  first-failure attribution.
\end{itemize}

\section{Problem Formulation and Evidence Boundary}
\subsection{Allowed Method Inputs}
\begin{itemize}
  \item Natural task language.
  \item Public RGB-D scene evidence and camera calibration.
  \item Public tool meshes and reviewed-or-pending tool-part evidence.
  \item Optional robot/gripper state needed for execution.
\end{itemize}

\subsection{Forbidden Method Inputs}
\begin{itemize}
  \item Hidden evaluator target refs, oracle target axes, oracle trajectories.
  \item Simulator target centers used to snap VLM refs.
  \item Tool-family or task-family fallback.
  \item Primitive fallback, repair-loop fallback, seed-specific shortcuts.
\end{itemize}

\subsection{Reference Representation}
\begin{align}
\mathcal{R}_{tool} &= \{p^{contact}_{tool}, a^{motion}_{tool}\}, \\
\mathcal{R}_{scene} &= \{p^{start}_{scene}, p^{path}_{scene},
p^{end}_{scene}, a^{direction}_{scene}\}.
\end{align}

\section{Method: ToolFrame}
\subsection{A. Tool-Centric Prior Extraction}
\begin{itemize}
  \item Input: tool mesh, P3SAM/SAMPart3D part evidence, task language.
  \item VLM selects a functional part from segmented evidence.
  \item Contact point and motion axis are selected only on the chosen part.
  \item Output: tool.contact\_point and tool.motion\_axis.
\end{itemize}

\subsection{B. Scene Key-Points Grounding}
\begin{itemize}
  \item Input: public RGB-D image, depth, camera calibration, task language.
  \item VLM marks start\_point, optional path\_points, end\_point, and
  scene.direction.
  \item RGB-D lifting converts 2D refs into 3D scene refs.
  \item No hidden simulator target centers or oracle axes are used.
\end{itemize}

\subsection{C. Geometry-Constrained Optimizer}
\begin{itemize}
  \item Inputs: lifted scene refs, tool.contact\_point, tool.motion\_axis.
  \item Stages: grasp/lift, align, approach, action.
  \item Constraint vocabulary: point\_match, point\_offset, axis\_align,
  move\_along, orientation\_offset\_about\_axis, clearance.
  \item Output: whole-tool $SE(3)$ subgoals and simtoolreal\_fixed\_goals.json.
\end{itemize}

\subsection{D. Policy Rollout and Evaluation}
\begin{itemize}
  \item SimToolReal fixed-goal policy executes compiled subgoals.
  \item Task contracts score final success after rollout.
  \item Failure attribution labels: functional grounding, scene grounding,
  subgoal validity, policy execution, task success, needs\_review.
\end{itemize}

\section{Task-Contract Benchmark}
\subsection{Task Tiers}
\begin{table}[h]
\centering
\caption{Task tiers for reporting.}
\begin{tabular}{p{0.20\columnwidth}p{0.68\columnwidth}}
\toprule
Tier & Tasks \\
\midrule
Headline & brush\_sweep/forward; eraser\_wipe/tabletop\_wipe \\
Candidate & brush\_sweep/left; hammer\_strike/down; screwdriver\_insert/wide\_socket; marker\_trace/write\_l,c \\
Diagnostic & screwdriver\_insert/deep\_socket; dense screwdriver variants \\
Draft & side-board eraser, hammer, screwdriver tasks \\
\bottomrule
\end{tabular}
\end{table}

\subsection{Task Contracts}
\begin{itemize}
  \item brush\_sweep/red\_zone: final target-block occupancy inside the visible
  red outline zone; directional displacement is diagnostic.
  \item eraser\_wipe/tabletop\_wipe: diagnostic path gate, lifted refs on visible HDU
  sticker, flat-eraser contact coverage $\geq 30\%$ of target long-axis extent;
  reviewed/current.
  \item hammer\_strike/down: provisional semantic contact and strike-direction
  gate.
  \item screwdriver\_insert/wide\_socket: provisional tip radial alignment and
  insertion-depth gate.
  \item marker\_trace/write\_l,c: provisional hard side-board task judged by
  produced marker-tip ink.
\end{itemize}

\subsection{Metrics}
\begin{table*}[t]
\centering
\caption{Headline result table template. Fill only from frozen
launcher\_summary files and matching per-task summary.json files.}
\scriptsize
\setlength{\tabcolsep}{3pt}
\begin{tabular}{p{0.17\textwidth}p{0.10\textwidth}p{0.22\textwidth}cccccc}
\toprule
Task & Status & Review note & N & Func. & Scene & Subgoal & Policy & Success \\
\midrule
brush\_sweep/forward & provisional & success predicate reviewed; red-brush labels pending & \todo{} & \todo{} & \todo{} & \todo{} & \todo{} & \todo{} \\
eraser\_wipe/tabletop\_wipe & reviewed & current contract; flat-eraser labels approved & \todo{} & \todo{} & \todo{} & \todo{} & \todo{} & \todo{} \\
\bottomrule
\end{tabular}
\end{table*}

\section{Experiments}
\subsection{Main Questions}
\begin{enumerate}
  \item Can public image-grounded evidence produce valid whole-tool subgoals
  without hidden evaluator refs?
  \item Which stage limits end-to-end task success?
  \item How do headline tasks differ from candidate and diagnostic tasks?
  \item What failures become attributable when VLM grounding, compilation,
  policy execution, and task-effect evaluation are separated?
\end{enumerate}

\subsection{Main Results}
\begin{itemize}
  \item Report headline task success rates only after final frozen rerun.
  \item Do not use reports/mainline\_6.4/ as final success statistics.
  \item Candidate and diagnostic tasks are reported separately as stress or
  readiness evidence.
\end{itemize}

\subsection{Ablations and Diagnostics}
\begin{table}[h]
\centering
\caption{Boundary-preserving ablations and diagnostics.}
\begin{tabular}{p{0.35\columnwidth}p{0.55\columnwidth}}
\toprule
Diagnostic & Purpose \\
\midrule
Public human-clicked scene refs & Upper-bound scene grounding diagnostic; no hidden evaluator refs. \\
Tool-reference sensitivity & Measure sensitivity to contact point and motion-axis choices. \\
Remove tool motion axis & Should fail preflight/compile; verifies no family-axis guessing. \\
Remove selected contact point & Tests need for explicit functional point. \\
No RGB-D lifting & Tests metric scene grounding dependence. \\
Strict tracking diagnostic & Measures tracking/effect disagreement; does not replace task contract. \\
\bottomrule
\end{tabular}
\end{table}

\section{Figures and Tables}
\begin{itemize}
  \item Fig. 1: Overview pipeline: tool refs, scene refs, compiler, rollout,
  evaluator.
  \item Fig. 2: Tool-reference evidence: selected part, contact point, motion
  axis.
  \item Fig. 3: Scene grounding overlays and lifted RGB-D refs.
  \item Fig. 4: Compiled whole-tool $SE(3)$ subgoals during rollout.
  \item Fig. 5: Failure attribution stacked bar by task.
  \item Table I: Task tiers and review status.
  \item Table II: Headline benchmark results.
  \item Table III: Candidate/diagnostic task readiness.
  \item Table IV: Ablations and diagnostics.
\end{itemize}

\section{Limitations}
\begin{itemize}
  \item Requires visible RGB-D evidence and reliable lifting.
  \item Requires public tool-part evidence and human review for final labels.
  \item Fixed-goal policy can fail even with valid compiled subgoals.
  \item Candidate tasks such as insertion and marker tracing need additional
  contract review before headline reporting.
  \item No online repair-loop fallback or seed-specific shortcuts.
\end{itemize}

\section{Conclusion}
\begin{itemize}
  \item Reiterate: VLMs select auditable evidence; the compiler solves
  whole-tool geometry; task contracts evaluate executed effects.
  \item Main takeaway after final rerun:
  \todo{insert final quantitative conclusion and main bottleneck.}
\end{itemize}

\section*{Appendix Checklist}
\begin{itemize}
  \item Full task contracts and thresholds.
  \item Tool-ref evidence manifests.
  \item Scene grounding prompts and overlay examples.
  \item Constraint program examples.
  \item Frozen launcher\_summary files and per-task summary.json files.
  \item Failure attribution examples.
  \item Asset and license statement.
\end{itemize}

\end{document}
